\documentclass[11pt]{article}
\usepackage[a4paper,margin=25mm]{geometry}
\usepackage{amsmath,amssymb}
\usepackage[T1]{fontenc}
\usepackage[hidelinks]{hyperref}
\setlength{\emergencystretch}{3em}
\title{The stated gamma expansion does not exist:\\a counterexample to Conjecture 7676}
\author{gaochengzhecpu (AI-assisted submission)}
\date{3 October 2026}
\begin{document}
\maketitle

\paragraph{The assertion.}
Conjecture \texttt{00000007676} explicitly defines
\[
 A_n(q,t)=\sum_{\pi\in S_n}q^{\operatorname{maj}(\pi)}t^{\operatorname{des}(\pi)}
\]
and asserts an expansion
\[
 A_n(q,t)=\sum_{j=0}^{\lfloor(n-1)/2\rfloor}
 \gamma_{n,j}(q)t^j(1+t)^{n-1-2j},
 \qquad\gamma_{n,j}(q)\in\mathbb N[q].
\]
It also specifies the number of monomials of these gamma coefficients.
The displayed range of $j$ is the standard gamma-basis range in degree
$n-1$; see \cite[Introduction]{LZ}. We use exactly the pair of statistics
and the weight $q^{\operatorname{maj}}t^{\operatorname{des}}$ written in
the problem.

\paragraph{Definitions and exhaustive computation.}
For $\pi=\pi_1\cdots\pi_n$, its descent set is
$D(\pi)=\{i:1\le i<n,\ \pi_i>\pi_{i+1}\}$.
The descent number is $|D(\pi)|$, and the major index is
$\sum_{i\in D(\pi)}i$. There are exactly two permutations of two labels:
\[
\begin{array}{c|c|c|c|c}
 \pi & D(\pi) & \operatorname{maj}(\pi)&\operatorname{des}(\pi)&
 q^{\operatorname{maj}(\pi)}t^{\operatorname{des}(\pi)}\\\hline
 12 & \varnothing &0&0&1\\
 21 & \{1\} &1&1&qt
\end{array}
\]
Consequently $A_2(q,t)=1+qt$.

\paragraph{Contradiction.}
At $n=2$, the only gamma index is $j=0$. The claimed identity becomes
\[
 1+qt=\gamma_{2,0}(q)(1+t).
\]
The coefficient of $t^0$ forces $\gamma_{2,0}(q)=1$, while the
coefficient of $t^1$ forces $\gamma_{2,0}(q)=q$. These are distinct
polynomials, so there is no such gamma expansion at all.

Equivalently, specialize $q=2$ and write $c=\gamma_{2,0}(2)$. Evaluating
at $t=0$ gives $c=1$; at $t=1$ it gives $3=2c$, a contradiction.
This argument permits arbitrary integer $c$, so imposing nonnegative
polynomial coefficients cannot repair the identity. It is therefore
incorrect to describe this example as having a negative gamma
coefficient: the displayed expansion itself does not exist. The
additional monomial-count assertion cannot hold for the nonexistent
coefficients.

\newpage
\paragraph{Formal verification.}
The accompanying Lean 4.19.0 project imports only \texttt{Std}.
The predicate \texttt{IsPermutation} means bijectivity on \texttt{Fin 2}.
The two-letter representation encodes every map. The theorem
\begin{center}\texttt{permutationWords\_exact}\end{center}
identifies the two enumerated words with precisely all bijections.
The representation is injective, and the enumeration has no duplicates.
Internal labels $0,1$ correspond to ordinary labels $1,2$; their order
comparisons agree. The one possible descent receives position $1$, as
required for the major index.

\texttt{qEulerian} sums the genuine major-index and descent weights over
this complete enumeration. \texttt{eulerianTwo\_formula} proves
$A_2(q,t)=1+qt$ for all integer evaluations. Nonnegative coefficient
polynomials in $q$ are represented by finite lists of natural numbers
and evaluated by Horner's rule. \texttt{gammaSum} uses the complete
standard gamma-index range, which reduces to $\{0\}$ at $n=2$.

\texttt{no\_nonnegative\_gamma\_expansion} rules out any family of
such coefficient polynomials satisfying the displayed identity.
The final theorem \texttt{conjecture7676\_counterexample} uses the
verified full permutation enumeration to negate the $n=2$ instance of
the universal assertion. Its formulation asks for identity at all
integer $q,t$; a polynomial identity necessarily has that property, so
failure at the two displayed evaluations suffices to rule it out.

The theorem assumes no gamma coefficients exist in advance. There are
no proof placeholders, added axioms, or native decision procedures.
The printed axiom audit uses only Lean's standard \texttt{propext} and
\texttt{Quot.sound}.

\begin{thebibliography}{9}
\bibitem{LZ}
Z. Lin and J. Zeng,
\emph{The $\gamma$-positivity of basic Eulerian polynomials via group actions},
J. Combin. Theory Ser. A \textbf{135} (2015), 112--129.
\href{https://doi.org/10.1016/j.jcta.2015.04.006}{doi:10.1016/j.jcta.2015.04.006}.
\end{thebibliography}
\end{document}
